% Fragmento para el reporte. Requiere \usepackage{amsmath,amssymb}.
\subsection{Función de prueba y descenso por gradiente}
Para $d\in\{2,3\}$, la función de Rosenbrock es
\[
f(\mathbf{x})=\sum_{i=1}^{d-1}\left[100(x_{i+1}-x_i^2)^2+(1-x_i)^2\right],
\qquad \mathbf{x}\in[-2.048,2.048]^d.
\]
Su mínimo global es $f(\mathbf{1})=0$. La forma del valle permite
estudiar el comportamiento del método en una geometría con curvaturas
desiguales; conocer el mínimo facilita medir el error sin usarlo para guiar
la búsqueda.

La condición inicial $\mathbf{x}^{(0)}$ se proporciona explícitamente
por consola o configuración. Si se omite, se muestrea uniformemente
en el dominio con la semilla indicada.

El código aproxima cada componente del gradiente por diferencias centrales:
\[
g_{k,j}=\frac{f(\mathbf{x}_k+h\mathbf{e}_j)-f(\mathbf{x}_k-h\mathbf{e}_j)}{2h},
\qquad h>0.
\]
La actualización es un paso proyectado:
\[
\mathbf{z}_k(\alpha)=\Pi_{\mathcal{D}}(\mathbf{x}_k-\alpha\mathbf{g}_k),
\qquad \mathcal{D}=[-2.048,2.048]^d.
\]
La búsqueda por retroceso comienza con $\alpha_0>0$ y reduce
$\alpha\leftarrow\rho\alpha$, con $0<\rho<1$, hasta satisfacer
la condición de Armijo proyectada implementada:
\[
f(\mathbf{z}_k(\alpha))\leq f(\mathbf{x}_k)
-c\,\mathbf{g}_k^{\mathsf T}(\mathbf{x}_k-\mathbf{z}_k(\alpha)),
\qquad 0<c<1.
\]
El punto aceptado define $\mathbf{x}_{k+1}$. El método se detiene
por tolerancia del gradiente, agotamiento del presupuesto,
falta de aceptación en el número máximo de retrocesos o ausencia
de desplazamiento. No se garantiza convergencia por fijar un presupuesto.

\subsection{Presupuesto de evaluación y reproducibilidad}
Un gradiente por diferencias centrales en $d$ dimensiones requiere
$2d$ evaluaciones de la función. Para evitar doble conteo:
\[
N_{\mathrm{equiv}}=N_f+2dN_{\nabla},
\]
donde $N_f$ cuenta las llamadas directas, incluyendo inicialización
y búsqueda lineal, y $N_{\nabla}$ cuenta las aproximaciones del gradiente.

En $S$ corridas se resumen media, desviación estándar, mejor, peor
y frecuencia de éxito según el umbral $\varepsilon$ predefinido:
\[
\overline{f}=\frac{1}{S}\sum_{s=1}^{S}f_s,\qquad
s_f=\sqrt{\frac{1}{S-1}\sum_{s=1}^{S}(f_s-\overline{f})^2},\qquad
\widehat{p}_{\mathrm{exito}}=\frac{1}{S}\sum_{s=1}^{S}\mathbb{I}(f_s<\varepsilon).
\]
Con GD y una condición inicial fija, variar la semilla no genera
trayectorias distintas en esta implementación. Para estudiar variabilidad
entre inicializaciones se conserva el muestreo del punto inicial o se
declara previamente un conjunto de condiciones distintas.
